% ============================================================================
%  SECTION 2 — RAPPELS ET CADRE THEORIQUE        (Membre 1)
% ============================================================================
\section{Rappels et cadre theorique}
\label{sec:rappels}

Cette section rassemble les notions du cours mobilisees dans la suite. Elle ne
vise pas l'exhaustivite mais fixe le vocabulaire, les notations et les
proprietes sur lesquelles reposent les sections suivantes.

\subsection{Probleme d'optimisation}

Un probleme d'optimisation consiste a minimiser (ou maximiser) une fonction
objectif $f : A \to \R$ sur un ensemble admissible $A$ :
\[
  \min_{x \in A} \; f(x).
\]
Un element $x \in A$ est une \emph{solution admissible} ; une solution qui
realise le minimum est une \emph{solution optimale}, notee $\xopt$. Lorsque
l'ensemble admissible est decrit par des contraintes d'inegalite, on ecrit
\[
  \min_{x \in \R^{n}} f(x)
  \quad\text{sous contraintes}\quad
  c_i(x) \le b_i, \quad i = 1, \dots, m.
\]

\subsection{Convexite}

\begin{definition}[Ensemble convexe]
Un ensemble $A \subseteq \R^{n}$ est \emph{convexe} si, pour tout couple de
points de $A$, le segment qui les joint est entierement contenu dans $A$ :
\[
  \forall x, y \in A,\ \forall \lambda \in [0,1] :\quad
  \lambda x + (1-\lambda) y \in A.
\]
\end{definition}

\begin{definition}[Fonction convexe]
Une fonction $f : \R^{n} \to \R$ est \emph{convexe} si
\[
  \forall x, y \in \R^{n},\ \forall \lambda \in [0,1] :\quad
  f\big(\lambda x + (1-\lambda) y\big)
  \le \lambda f(x) + (1-\lambda) f(y).
\]
Geometriquement, tout segment reliant deux points du graphe se situe au-dessus
(ou sur) le graphe. Lorsque l'inegalite est \emph{stricte} pour $x \neq y$ et
$\lambda \in {]0,1[}$, la fonction est dite \emph{strictement convexe}.
\end{definition}

Deux proprietes de stabilite, etablies en travaux diriges, seront utilisees
pour montrer que notre probleme est convexe :
\begin{proposition}[Stabilite de la convexite]
\label{prop:stabilite}
\leavevmode
\begin{enumerate}
  \item Un demi-espace $\{x \in \R^{n} : a\T x \le b\}$ est un ensemble convexe ;
        une intersection d'ensembles convexes est convexe.
  \item Une somme de fonctions convexes est convexe ; le produit d'une fonction
        convexe par un scalaire positif est convexe.
\end{enumerate}
\end{proposition}

\subsection{La propriete fondamentale, et ses limites}

\begin{theoreme}[Propriete-or]
\label{th:or}
Tout minimum local d'une fonction convexe est un minimum global.
\end{theoreme}

C'est cette propriete qui justifie l'interet des problemes convexes : un
algorithme qui converge vers un minimum local converge en realite vers
l'optimum global. Une precision s'impose toutefois, car elle conditionne la
lecture de nos resultats experimentaux.

\begin{remarque}[Convexite et unicite]
\label{rem:unicite}
La convexite garantit l'unicite de la \emph{valeur optimale} $\Fopt$, mais
\textbf{pas} celle du \emph{minimiseur} $\xopt$. L'ensemble des minimiseurs
peut etre reduit a un point, mais aussi former tout un segment --- le cours en
donne d'ailleurs un exemple en programmation lineaire, ou l'ensemble des
solutions optimales est un segment entier. L'unicite du minimiseur requiert une
hypothese plus forte, la \emph{stricte} convexite. Comme nous le verrons
(section~\ref{sec:probleme}), le LASSO n'est pas strictement convexe en general ;
on ne pourra donc affirmer que la convergence se fait vers une \emph{meme valeur}
$\Fopt$, et non necessairement vers un \emph{meme point} $\xopt$.
\end{remarque}

\subsection{Programmation lineaire et quadratique}

Le cours a etudie en detail la \emph{programmation lineaire}, cas particulier ou
objectif et contraintes sont lineaires, resolue par l'algorithme du simplexe.
Notre probleme generalise ce cadre dans une direction differente : l'objectif y
sera \emph{quadratique} (donc convexe mais non lineaire) et augmente d'un terme
de regularisation \emph{non differentiable}. Ni le simplexe ni la descente de
gradient classique ne s'y appliquent directement, ce qui motive l'introduction
des outils de la section~\ref{sec:theorie}.
